\textbf{Proof.}
Assumption \(\texttt{ass:coordinate-bounded-kappa-feasibility}\) makes \(\mathcal B_{\kappa,\overline a}\) nonempty.  Put \(N=2n\), index the exposure coordinates by \(\ell\), and write \(Z_\ell=D_\ell-\pi_\ell\).  Lemma \(\texttt{lem:directed-local-centered-moment-count}\), whose proof partitions all five fourth-order multiplicity types and the directed one-cycle, two-cycle, inward-star, chain, three-cycle, and four-cycle cases, gives
\[
\mathfrak M_{3,n}\le8N d_n^2,
\qquad
\mathfrak M_{4,n}\le64N^2d_n^2,
\qquad
s_n\le d_n\le\overline d n^\alpha.
\tag{1}
\]
These are consequences only of Zhao's coordinate-versus-complement local independence condition.

Fix a bounded exposure-potential-outcome array and \(B\in\mathcal B_{\kappa,\overline a}\).  On \(D_{i,a}=1\), known exposure gives \(Y_i(W)=Y_i(a)\).  With the signed convention in \(Y\),
\[
\mathbb E_\nu\widehat Y_\ell^{\mathrm{obs}}=Y_\ell,
\qquad
\{\Pi^{-1}(\widehat Y^{\mathrm{obs}}-Y)\}_\ell
=\frac{Y_\ell Z_\ell}{\pi_\ell^2}.
\tag{2}
\]
Thus there is no exposure-mismatch term.  The pseudoinverse is the ordinary inverse on this rank-two class.  Consequently, with
\[
G_B=\widehat\gamma^{\mathrm{HT}}(B)-\gamma_B^*(Y),
\]
equation (2) gives
\[
G_B=\sum_{\ell=1}^{N}u_{B,\ell}Z_\ell,
\qquad
u_{B,\ell}:=S_B^{-1}B^\top\Omega e_\ell\frac{Y_\ell}{\pi_\ell^2}.
\tag{3}
\]
This is the only coefficient name used below.  Every covariance entry has absolute value at most one, every covariance column has at most \(s_n\) nonzero entries, \(\lVert S_B^{-1}\rVert_{\mathrm{op}}\le(n\kappa)^{-1}\), and every basis coordinate is bounded by \(\overline a\).  Hence exposure-probability control and (1) give a finite \(K_u\), independent of \(n,B,Y,H\), such that
\[
\sup_{B,\ell}\lVert u_{B,\ell}\rVert_2
\le K_u n^{\alpha+2\beta-1}.
\tag{4}
\]
Also
\[
X_B=n^{-1}B^\top(D-\mathbb E_\nu D)
=\sum_{\ell=1}^{N}x_{B,\ell}Z_\ell,
\qquad
\sup_{B,\ell}\lVert x_{B,\ell}\rVert_2\le\overline a/n.
\tag{5}
\]
Applying the fourth-moment envelope in (1) coordinatewise to (3) and (5), and using \((z_1^2+z_2^2)^2\le2z_1^4+2z_2^4\), gives constants \(K_G,K_X<\infty\) for which
\[
\sup_B\mathbb E_\nu\lVert G_B\rVert_2^4
\le K_Gn^{6\alpha+8\beta-2},
\qquad
\sup_B\mathbb E_\nu\lVert X_B\rVert_2^4
\le K_Xn^{2\alpha-2}.
\tag{6}
\]
The feasible and oracle estimators obey
\[
\widehat\tau_B^{\mathrm{CV}}=\widehat\tau_B^{\mathrm{or}}+\Delta_B,
\qquad \Delta_B=-G_B^\top X_B.
\tag{7}
\]
Cauchy--Schwarz and (6) yield a finite \(K_\Delta\) such that
\[
\sup_B\mathbb E_\nu\Delta_B^2
\le K_\Delta n^{4\alpha+4\beta-2}.
\tag{8}
\]

Put \(T_B=\widehat\tau_B^{\mathrm{or}}-\tau\).  Equation (2) makes the Horvitz--Thompson part unbiased, and the basis statistic is centered, so \(\mathbb E_\nu T_B=0\).  Direct expansion gives
\[
T_B=\sum_{\ell=1}^{N}t_{B,\ell}Z_\ell,
\qquad
t_{B,\ell}=n^{-1}\{Y_\ell/\pi_\ell-B_{\ell\cdot}\gamma_B^*(Y)\}.
\tag{9}
\]
The oracle coefficient is an orthogonal projection in the \(\Omega^{1/2}\)-geometry, so
\[
\gamma_B^*(Y)^\top S_B\gamma_B^*(Y)
\le Y^\top\Pi^{-1}\Omega\Pi^{-1}Y.
\tag{10}
\]
There are at most \(Ns_n\) nonzero covariance entries.  Bounded outcomes, exposure positivity, and \(S_B\succeq n\kappa I_2\) therefore give finite constants \(K_\gamma,K_t\) such that
\[
\sup_B\lVert\gamma_B^*(Y)\rVert_2
\le K_\gamma n^{\alpha/2+\beta},
\qquad
\sup_{B,\ell}|t_{B,\ell}|
\le K_tn^{\alpha/2+\beta-1}.
\tag{11}
\]
Expanding \(\operatorname{Cov}_\nu(T_B,\Delta_B\mid Y)\) in the two coefficient coordinates and applying the third-moment envelope in (1) to the three linear forms in (3), (5), and (9) gives a finite \(K_{\mathrm{cov}}\) such that
\[
\sup_B|\operatorname{Cov}_\nu(T_B,\Delta_B\mid Y)|
\le K_{\mathrm{cov}}n^{(7/2)\alpha+3\beta-2}.
\tag{12}
\]
The exponent is the sum of the factors \(Nd_n^2\), \(\sup|t_{B,\ell}|\), \(\sup\lVert u_{B,\ell}\rVert_2\), and \(\sup\lVert x_{B,\ell}\rVert_2\).

The variance identity
\[
\operatorname{Var}_\nu(T_B+\Delta_B\mid Y)-\operatorname{Var}_\nu(T_B\mid Y)
=\operatorname{Var}_\nu(\Delta_B\mid Y)+2\operatorname{Cov}_\nu(T_B,\Delta_B\mid Y)
\tag{13}
\]
together with (8)--(12) proves the first uniform bound with
\[
R_n=K_\Delta n^{4\alpha+4\beta-2}
+2K_{\mathrm{cov}}n^{(7/2)\alpha+3\beta-2}.
\tag{14}
\]
Because \(\alpha,\beta\ge0\), \((7/2)\alpha+3\beta\le4\alpha+4\beta\).  Thus, with \(K_R=K_\Delta+2K_{\mathrm{cov}}\),
\[
0\le R_n\le K_Rn^{4\alpha+4\beta-2}.
\tag{15}
\]
The constants may depend on the fixed \(\kappa,\overline a,\underline\pi,\overline d\), but not on \(n,B,Y,H\).

Completing the square for the fixed oracle coefficient gives
\[
\operatorname{Var}_\nu(\widehat\tau_B^{\mathrm{or}}\mid Y)
=n^{-2}Y^\top\{\Pi^{-1}\Omega\Pi^{-1}
-\Pi^{-1}\Omega BS_B^{-1}B^\top\Omega\Pi^{-1}\}Y.
\tag{16}
\]
Under the i.i.d. potential-pair assumption,
\[
\mathbb E_{H^{\otimes n}}(YY^\top)=Q(c_H).
\tag{17}
\]
Taking the trace expectation in (16), using (17), and applying \(\texttt{def:gain-loss}\) proves
\[
\mathbb E_{H^{\otimes n}}
\operatorname{Var}_\nu(\widehat\tau_B^{\mathrm{or}}\mid Y)
=h(c_H)-g_B(c_H)=\ell_B(c_H).
\tag{18}
\]
Integrating (13)--(15) over \(H^{\otimes n}\) and applying (18) proves the second uniform bound.  Finally,
\[
nR_n\le K_Rn^{4\alpha+4\beta-1}\longrightarrow0
\]
by \(\texttt{ass:zhao-first-order-rate}\).  Every inner variance is conditional design variance under \(\nu\); the expectation under \(H^{\otimes n}\) is only the declared ex ante basis-performance average. \(\square\)